\documentclass{article}
\usepackage{/globals/code}
\usepackage{/globals/anki}
\ankishow
\ankiprefix{schleifen}
\ankitopic{Schleifen}
\begin{document}
\section{Schleifen} 
Schleifen werden genutzt um bei Bedarf Abschnitte eines Programms zu wiederholen solange eine bestimmte Bedingung erfüllt ist.

\subsection{\texttt{while}}
Eine \ilsourcecode{while}-Schleife wiederholt bestimmte Anweisungen solange eine gegebene Bedingung \ilsourcecode{true} ist.
\begin{sourcecode}{java}
while (bedingung) {
    Anweisungen;
}
\end{sourcecode}
Sie ist eine abweisende Schleife; wenn die Bedingung bei dem ersten Durchgang bereits \ilsourcecode{false} ist werden die Anweisungen kein einziges mal durchgeführt.
\ankinote{syntaxe}{Syntaxe}{while, do while, for}
\ankinote{arten}{Arten}{abweisend (while), annehmend (do)}

\subsection{\texttt{do while}}
Eine \ilsourcecode{do while}-Schleife ist der \ilsourcecode{while}-Schleife sehr ähnlich.
\begin{sourcecode}{java}
do {
    Anweisungen;
} while (bedingung)
\end{sourcecode}
Sie ist aber annehmend. Auch wenn die Bedingung von Anfang an \ilsourcecode{false} ist werden die Anweisungen mindestens ein mal ausgeführt.

\subsection{\texttt{for}}
Eine \ilsourcecode{for}-Schleife ist letztendlich eine \ilsourcecode{while}-Schleife; sie wird wiederholt solange die Bedingung \ilsourcecode{true} ist, aber sie bietet extra Syntax um eine Anweisung beim ersten Schleifenbeginn und eine Anweisung nach jedem Durchlauf auszuführen.
\begin{sourcecode}{java}
for(Initialisierung; bedingung; Aktualisierung) {
    Anweisungen;
}
\end{sourcecode}
So kann sehr einfach hochgezählt werden
\begin{sourcecode}{java}
for(int i = 0; i < 10; i++) {
    Anweisungen;
}
\end{sourcecode}
\ankinote{for}{for}{Initialisierung; bedingung; Aktualisierung, Aktualisierung nach jedem Durchgang}

\subsection{Äquivalenzen}
Alle Schleifenarten machen letztendlich das Gleiche. Jede \ilsourcecode{while} kann auch als \ilsourcecode{do while} geschrieben werden; wenn vorher bereits bekannt ist dass in jedem Fall die Bedingung mindestens beim ersten mal \ilsourcecode{true} sein wird, wird das \ilsourcecode{if} garnicht mal benötigt.
\begin{sourcecode}{java}
while (bedingung) {
    Anweisungen;
}
// ist äquivalent zu
if (bedingung) {
    do {
        Anweisungen;
    } while (bedingung)
}
\end{sourcecode}
Andersherum können \ilsourcecode{do while}s auch als \ilsourcecode{while}s geschrieben werden. 
\begin{sourcecode}{java}
do {
    Anweisungen;
} while (bedingung)
// ist äquivalent zu
Anweisungen;
while (bedingung) {
    Anweisungen;
}
\end{sourcecode}
Ein \ilsourcecode{for} ist nur eine kompakteres \ilsourcecode{while}.
\begin{sourcecode}{java}
for(Initialisierung; bedingung; Aktualisierung) {
    Anweisungen;
}
// ist äquivalent zu
{
    Initialisierung; 
    while(bedingung) {
        Anweisungen;
        Aktualisierung;
    }
}
\end{sourcecode}

\subsection{\texttt{break;} und \texttt{continue;}}
Ein \ilsourcecode{break;} bricht die zurzeitige Schleife ab, auch wenn die Bedingung noch \ilsourcecode{true} sein mag.

Mit einem \ilsourcecode{continue;} wird der Rest des zurzeitigen Durchgangs übersprungen, die Schleife kann aber noch weiter laufen, wenn die Bedingung \ilsourcecode{true} ist.

\subsection{Endlosschleifen}
Eine Schleife die kein Ende hat und sich für immer weiter und weiter wiederholt ist eine Endlosschleife. Sie terminiert nicht.
\begin{sourcecode}{java}
while (true) {
   ;
}
\end{sourcecode}
\ankinote{endlos}{Endlosschleife}{Terminiert nicht}
\end{document}